\section{The Substrate as a Computer}

A world that is a single rule stepped over a graph is, whatever else it is, a computer, and it is worth asking what kind. The answer is a powerful one, and it is shaped in a particular way by the negative curvature, which makes the substrate not merely universal but unusually well suited to the kind of associative, content-addressed computation that minds appear to do.

That the knit can compute anything follows on several independent legs, of increasing strength. The local rule is functionally complete, able to embed a universal logic gate and the elementary rule that is known to be Turing complete on its own. A register machine, with its counters and tests, can be laid out so that its operations respect charge conservation, giving universality with the base's own bookkeeping intact. And the strongest leg pins universality to the geometry, by realizing Conway's Game of Life directly on the flat cubic cusp, with no background scaffolding the rule does not provide. Weak universality would tolerate a fixed helping structure, strong universality does not, and the substrate reaches the strong kind on the very space that is physical three-dimensional space.

\subsection{The bulk is a database}

The more striking computational fact is about layout rather than logic. Because the bulk branches like a tree, with a fan-out of twenty-four and a capacity that grows like $e^{3R}$ with radius, it has logarithmic diameter, every dock reachable from every other in a number of steps that grows only as the logarithm of the size. A structure that costs linear or polynomial depth on a flat array costs logarithmic depth here. Addresses make this concrete. Every dock has a unique, self-describing address, a word in the reflection group, and from that word alone its neighbours can be reconstructed by a finite automaton and a route to any target found by pure address arithmetic, greedily, following at each step the neighbour that lies nearest the destination. On the committed substrate this greedy routing succeeds every time with a path barely longer than the true shortest one, and the core adjacency holds out to millions of docks. Tree-shaped data structures therefore instantiate in the bulk with no stored pointers at all, the geometry supplying the links that a flat machine would have to keep in memory.

\subsection{Memory by content}

What a mind does with memory is not look up an address, it recalls a whole from a fragment, and that is exactly the operation the bulk is built for. Treat each dock as holding a stored word and the recall is content-addressable, a cue broadcast to every dock at once, each deciding locally whether it matches, with no central scan. This is the associative-computing paradigm, and the substrate realizes it with the geometry doing the heavy lifting \cite{potter1992}. The broadcast itself, a wave spreading outward shell by shell, is what cognitive science calls spreading activation, and because the bulk has logarithmic diameter the wave reaches the entire memory in a handful of beats. The result is that search latency scales logarithmically with the size of the memory rather than as a polynomial, and capacity within a ball of fixed radius grows exponentially rather than polynomially, both measured against a flat cubic control that comes out far worse on each (depth L3).

\begin{table}[ht]
\centering
\begin{tabular}{@{}lll@{}}
\toprule
quantity & on $\{3,4,3,4\}$ & on a flat cubic memory \\
\midrule
search latency with size & logarithmic & polynomial, like $N^{1/3}$ \\
capacity per unit radius & exponential & polynomial \\
parallel match cost & constant in size & constant in size \\
\bottomrule
\end{tabular}
\caption{The associative memory on the curved bulk against a flat control. The curvature, not the bit budget, is the advantage.}
\label{tab:assoc}
\end{table}

Run across the whole family of tessellations, this advantage tracks curvature directly, the more negatively curved the substrate the higher the capacity and the lower the latency, with the flat lattice the worst on both. The classic data structures, hash tables, balanced trees, priority queues, all realize on the bulk with the same logarithmic depth, which is why the later chapters can read the bulk as the storehouse of a mind without inventing any new mechanism. The mental realm, when we reach it, will simply be this database read from the inside.
